\section{The deterministic argument also uses path norm directly}
The energy bound in the theorem is a convenient sufficient condition.
The actual projection argument has a more general form. Put
\[
 P=\frac1N\sum_i|a_i|\|U_i\|,\qquad
 A_2=\frac1N\sum_i a_i^2,\qquad U_2=\frac1N\sum_i\|U_i\|^2.
\]
Suppose instead that $P\le(1+\varepsilon)s/\rho$ and
$\E(s-\chi F)_+\le h$, without a bound on the balanced quadratic energy.
Then the same comparator $q_A$ obeys (4), with
\[
 u_0=\sqrt{(s+P\sqrt D)(\varepsilon s+2h)},\qquad
 v_0=\sqrt{2P(\varepsilon s+h)/\rho},
\]
and its norm and feature radius satisfy
\[
 \|q_A\|_{\mathcal H(K_e)}^2
       \le A_2\left(\frac{\rho}{s b_j}\right)^2,
 \qquad \max_xK_e(x,x)\le D U_2.                    \tag{11a}
\]
Indeed (9)--(10) already use $P$ before replacing it by $E$.
Also $|F|\le P\sqrt D$, and the unchanged reweighting uses the original
readouts, whose mean square is $A_2$. No rescaling of the frozen features
is needed. This form is useful for non-Gaussian priors: a near-minimum
power cost can sometimes imply a near-minimum path norm even when its
quadratic parameter energy is larger.
